\section{Next steps}

\begin{frame}{Limitations}
  \small
  \begin{enumerate}
    \item \alert{No timestep resampling \& discretization:} cases are not yet
          projected onto a 15/30/60-min grid or binned into Low/Medium/High, so
          $X_T$ for the IL trigger is not built.
    \item \alert{Appliance-only detection:} the engine supports only
          \texttt{power\_cycle}; the comfort inquiry produced 0 cases.
    \item \alert{No cooperative quality control:} users cannot confirm, edit, or
          reject bot-inferred labels.
    \item \alert{No occupancy sensors:} without PIR or door contacts,
          \texttt{who} depends on user input.
    \item \alert{Abnormally long fridge cycles} (3 of 48) are not yet split or
          explained.
  \end{enumerate}
\end{frame}

\begin{frame}{Future work}
  \footnotesize
  \begin{columns}[T,onlytextwidth]
    \begin{column}{0.47\textwidth}
      \begin{block}{Timestep discretization}
        Resample cases onto a standard grid (15/30/60 min) to build $X_T$, then
        bin values into Low/Medium/High for the IL trigger.
      \end{block}
      \begin{block}{Climate binding}
        Attach temperature and humidity to each TV interval
        $[t_{\mathrm{start}}, t_{\mathrm{end}}]$ and add comfort-band rules.
      \end{block}
    \end{column}
    \begin{column}{0.47\textwidth}
      \begin{block}{Cooperative QC}
        Web UI to Confirm / Edit / Reject / Defer bot-filled labels.
      \end{block}
      \begin{block}{Occupancy sensing}
        PIR, door contact, and CO$_2$ to estimate headcount (\texttt{who}).
      \end{block}
    \end{column}
  \end{columns}
\end{frame}

\usetikzlibrary{positioning,arrows.meta}

\begin{frame}{Pipeline and roadmap}
  \centering
  \begin{tikzpicture}[
      stage/.style={draw=templateRed, fill=templateRed!7, rounded corners=3pt,
        minimum width=2.3cm, minimum height=1.25cm, align=center,
        font=\footnotesize},
      todo/.style={stage, fill=templateDarkRed, text=white},
      arr/.style={-{Stealth}, thick, templateRed}]
    \node[stage] (a) {Raw sensor\\data};
    \node[stage, right=0.35cm of a] (b) {Cycle detection\\(adaptive thresholds)};
    \node[stage, right=0.35cm of b] (c) {Case\\$[t_{\mathrm{start}}, t_{\mathrm{end}}]$};
    \node[todo,  right=0.35cm of c] (d) {Timestep resampling\\\& discretization};
    \node[stage, right=0.35cm of d] (e) {IL trigger\\+ QC};
    \draw[arr] (a)--(b); \draw[arr] (b)--(c); \draw[arr] (c)--(d); \draw[arr] (d)--(e);
    \node[below=0.25cm of d, font=\scriptsize, templateDarkRed] {next milestone};
    \node[below=0.25cm of e, font=\scriptsize] {in progress};
    \node[below=0.25cm of b, font=\scriptsize] {done};
  \end{tikzpicture}

  \vspace{0.1cm}
  \footnotesize
  \raggedright
  \begin{itemize}
    \setlength{\itemsep}{1pt}
    \item \alert{Next step, after a case exists:} \emph{temporal resampling}
          projects cases of different lengths onto a standard time grid
          (15/30/60 min) to build the analysis matrix $X_T$; \emph{semantic
          binning} turns numeric values into Low / Medium / High.
    \item \alert{Why:} the IL trigger (neighbourhood density) needs time-aligned
          feature vectors to compare cases and decide whether to ask the user.
  \end{itemize}
\end{frame}

\begin{frame}{References}
  \footnotesize
  \begin{thebibliography}{9}
    \setbeamertemplate{bibliography item}[text]
    \bibitem{otsu} N.~Otsu, ``A threshold selection method from gray-level histograms,'' \emph{IEEE Trans. Syst., Man, Cybern.}, vol.~9, no.~1, pp.~62--66, 1979.
    \bibitem{neuralnilm} J.~Kelly and W.~Knottenbelt, ``Neural NILM: Deep neural networks applied to energy disaggregation,'' in \emph{Proc. ACM BuildSys}, 2015, pp.~55--64.
    \bibitem{precioso} D.~Precioso and D.~G\'omez-Ullate, ``Thresholding methods in non-intrusive load monitoring,'' \emph{J. Supercomputing}, vol.~79, pp.~14039--14062, 2023.
    \bibitem{jumiawi} W.~A.~H.~Jumiawi and A.~El-Zaart, ``Improvement in the between-class variance based on lognormal distribution for accurate image segmentation,'' \emph{Entropy}, 2022.
    \bibitem{ukdale} J.~Kelly and W.~Knottenbelt, ``The UK-DALE dataset, domestic appliance-level electricity demand and whole-house demand from five UK homes,'' \emph{Scientific Data}, vol.~2, 150007, 2015.
  \end{thebibliography}
\end{frame}
